%% ****** Start of file apstemplate.tex ****** %
%%
%%
%%   This file is part of the APS files in the REVTeX 4.2 distribution.
%%   Version 4.2a of REVTeX, January, 2015
%%
%%
%%   Copyright (c) 2015 The American Physical Society.
%%
%%   See the REVTeX 4 README file for restrictions and more information.
%%
%
% This is a template for producing manuscripts for use with REVTEX 4.2
% Copy this file to another name and then work on that file.
% That way, you always have this original template file to use.
%
% Group addresses by affiliation; use superscriptaddress for long
% author lists, or if there are many overlapping affiliations.
% For Phys. Rev. appearance, change preprint to twocolumn.
% Choose pra, prb, prc, prd, pre, prl, prstab, prstper, or rmp for journal
%  Add 'draft' option to mark overfull boxes with black boxes
%  Add 'showkeys' option to make keywords appear
%\documentclass[aps,prl,preprint,groupedaddress]{revtex4-2}
\documentclass[aps,prl,twocolumn,groupedaddress]{revtex4-2}
%\documentclass[aps,prl,preprint,superscriptaddress]{revtex4-2}
%\documentclass[aps,prl,reprint,groupedaddress]{revtex4-2}

% You should use BibTeX and apsrev.bst for references
% Choosing a journal automatically selects the correct APS
% BibTeX style file (bst file), so only uncomment the line
% below if necessary.
%\bibliographystyle{apsrev4-2}

\usepackage{graphicx}
\usepackage{epstopdf}
\usepackage{caption}
%\usepackage{amsmath}% http://ctan.org/pkg/amsmath
%\usepackage{amsthm}
%\usepackage{amsfonts}
%\usepackage{subfigure}
%\usepackage{hhline}
%\usepackage[miktex]{gnuplottex}
%\usepackage{xcolor}
\usepackage{amssymb}
\usepackage{amsmath}
\usepackage{color}
\usepackage{hyperref}
%\usepackage[percent]{overpic}
\usepackage{tikz}
\usepackage{mathrsfs}
\usepackage{wasysym}
\usepackage{tikz-cd}
\usepackage{linearb}
%\usepackage{wsuipa}
\usepackage{calc}
\usepackage{tipa}

%\usepackage{stix} %\fisheye
\usepackage{stackengine,scalerel}

% so sections, subsections, etc. become numerated.
\setcounter{secnumdepth}{3}

\newcommand{\FF}{\mathbb{F}}
\newcommand{\NN}{\mathbb{N}}
\newcommand{\ZZ}{\mathbb{Z}}
\newcommand{\QQ}{\mathbb{Q}}
\newcommand{\RR}{\mathbb{R}}
\newcommand{\CC}{\mathbb{C}}
\newcommand{\II}{\mathbb{I}}
\newcommand{\GG}{\mathbb{G}}

\DeclareMathOperator*{\argmax}{arg\,max}
\DeclareMathOperator*{\argmin}{arg\,min}
\newcommand{\avrg}[1]{\left\langle #1 \right\rangle}
\newcommand{\nelta}{\bar{\delta}}
\newcommand{\bra}[1]{\left\langle #1\right|}
\newcommand{\ket}[1]{\left| #1 \right\rangle}
\newcommand{\sbra}[1]{\langle #1|}
\newcommand{\sket}[1]{| #1 \rangle}
\newcommand{\bek}[3]{\left\langle #1 \right| #2 \left| #3 \right\rangle}
\newcommand{\sbek}[3]{\langle #1 | #2 | #3 \rangle}
\newcommand{\braket}[2]{\left\langle #1 \middle| #2 \right\rangle}
\newcommand{\ketbra}[2]{\left| #1 \middle\rangle \middle\langle #2  \right|}
\newcommand{\sbraket}[2]{\langle #1 | #2 \rangle}
\newcommand{\sketbra}[2]{| #1 \rangle  \langle #2 |}
\newcommand{\norm}[1]{\left\lVert#1\right\rVert}
\newcommand{\snorm}[1]{\lVert#1\rVert}
\newcommand{\bvec}[1]{\boldsymbol{\mathsf{#1}}}
\newcommand{\bcov}[1]{\boldsymbol{#1}}
\newcommand{\bdua}[1]{\boldsymbol{\check{#1}}}
%\newcommand{\bdov}[1]{\boldsymbol{\breve{#1}}}
\newcommand{\bdov}[1]{\breve{#1}}
%\newcommand{\bten}[1]{\boldsymbol{\mathfrak{#1}}}
\newcommand{\bten}[1]{\boldsymbol{\mathfrak{#1}}}
\newcommand{\forany}{\tilde{\forall}}
\newcommand{\qed}{$\overset{\circ}{.}\;$}
\newcommand{\bigo}{\mathcal{O}}
\newcommand{\spn}{\mathrm{spn}\,}
\newcommand{\krn}{\mathrm{krn}\,}
\newcommand{\biy}{\leftrightarrow}
\newcommand{\ugh}{\mathbin{\textlinb{\BPwheel}}}

\newcommand\bigeye{\ensurestackMath{\stackinset{c}{}{c}{-.3pt}%
  {\bullet}{\scriptstyle\bigcirc}}}
\newcommand\eye{\scalerel*{\bigeye}{x}}
%\newcommand*{\fisheye}{%
%    \mathbin{%
%        \ooalign{$\circledcirc$\cr\hidewidth$\bullet$\hidewidth}%
%    }%
%}

% double angle-bracket
% https://tex.stackexchange.com/questions/79657/how-to-get-double-angle-bracket-without-using-mnsymbol-package
\makeatletter
\newsavebox{\@brx}
\newcommand{\llangle}[1][]{\savebox{\@brx}{\(\m@th{#1\langle}\)}%
  \mathopen{\copy\@brx\mkern2mu\kern-0.9\wd\@brx\usebox{\@brx}}}
\newcommand{\rrangle}[1][]{\savebox{\@brx}{\(\m@th{#1\rangle}\)}%
  \mathclose{\copy\@brx\mkern2mu\kern-0.9\wd\@brx\usebox{\@brx}}}
\makeatother

% https://tex.stackexchange.com/questions/297362/bar-through-letter-in-mathmode
%\usepackage{calc}
\newsavebox\CBox
\newcommand\hcancel[2][0.5pt]{%
  \ifmmode\sbox\CBox{$#2$}\else\sbox\CBox{#2}\fi%
  \makebox[0pt][l]{\usebox\CBox}%  
  \rule[0.77\ht\CBox-#1/2]{\wd\CBox}{#1}}
\newcommand{\cb}{\hcancel{b}}
\newcommand{\cd}{\hcancel{d}}
\newcommand{\ct}{\hcancel{\partial}}

\begin{document}

% Use the \preprint command to place your local institutional report
% number in the upper righthand corner of the title page in preprint mode.
% Multiple \preprint commands are allowed.
% Use the 'preprintnumbers' class option to override journal defaults
% to display numbers if necessary
%\preprint{}

%Title of paper
\title{
Title...
}

% repeat the \author .. \affiliation  etc. as needed
% \email, \thanks, \homepage, \altaffiliation all apply to the current
% author. Explanatory text should go in the []'s, actual e-mail
% address or url should go in the {}'s for \email and \homepage.
% Please use the appropriate macro foreach each type of information

% \affiliation command applies to all authors since the last
% \affiliation command. The \affiliation command should follow the
% other information
% \affiliation can be followed by \email, \homepage, \thanks as well.
%\author{Juan I. Perotti}
%\email[]{juan.perotti@unc.edu.ar}
%\homepage[]{Your web page}
%\thanks{}
%\altaffiliation{}
%\affiliation{}
%\affiliation{Instituto de F\'isica Enrique Gaviola (IFEG-CONICET), Ciudad Universitaria, 5000 C\'ordoba, Argentina}
\affiliation{Facultad de Matem\'atica, Astronom\'ia, F\'isica y Computaci\'on, Universidad Nacional de C\'ordoba, Ciudad Universitaria, 5000 C\'ordoba, Argentina}

% \author{Colaborador}
% \affiliation{Instituto de F\'isica Enrique Gaviola (IFEG-CONICET), Ciudad Universitaria, 5000 C\'ordoba, Argentina}
% \affiliation{Facultad de Matem\'atica, Astronom\'ia, F\'isica y Computaci\'on, Universidad Nacional de C\'ordoba, Ciudad Universitaria, 5000 C\'ordoba, Argentina}
% \email[E-mail: ]{xxx.yyy@unc.edu.ar }

%Collaboration name if desired (requires use of superscriptaddress
%option in \documentclass). \noaffiliation is required (may also be
%used with the \author command).
%\collaboration can be followed by \email, \homepage, \thanks as well.
%\collaboration{Claudio J. Tessone}
%\noaffiliation

\date{\today}

%\begin{abstract}
%Bla bla bla... resumen...
%\end{abstract}

% insert suggested keywords - APS authors don't need to do this
%\keywords{}

%\maketitle must follow title, authors, abstract, and keywords
\maketitle
\onecolumngrid

\section{Theoretical review}

\subsection{Simplicial complex}

Consider labels $i,j,k,... \in \{1,2,...,N\}$.
A simplicial complex is a set $K$ of nodes identified by labels $i$, links identified by labels $ij$ with $i<j$, triangles identified by labels $ijk$ with $i<j<k$, and so on.
More generally, the elements of $K$ are sequences $i_0i_1...i_k$ called $k$-simplices that satisfy $0\leq i_0<i_1<...<i_k\leq N$.
In particular, a simplicial complex $K$ satisfies the condition of topological closure, meaning that if $i_0...i_k$ is a simplex in $K$, then $i_1...i_{j-1}i_{j+1}...i_k$ is also a simplex of $K$ for any $j=0,1,...,k$.

\subsection{Chains}

To each $k$-simplex $i_0...i_k$ of a simplicial complex $K$ we associate a vector $e_{i_0...i_k}$.
The vector space $C$ generated by the linear combinations of the set $\{e_{i_0...i_k}:i_0...i_k\in K\}$ is called the space of chains.
The set $\{e_{i_0...i_k}:i_0...i_k\in K\}$ denotes the canonical basis of $C$.
Any chain $c\in C$ can be written as
$$
c = 
\sum_{k=0}^N 
\sum_{i_0...i_k \in K}
c^{i_0...i_k}
e_{i_0...i_k}
$$
where the $c^{i_0...i_k}$ are just scalar coefficients denoting the coordinates of $c$ in the canonical basis.

\subsection{Cochains}

The dual space $C^*$ if $C$ is called the space of cochains.
The elements of the dual basis is are denoted by
$$
e^{i_0...i_k}
$$
and satisfy
$$
e^{i_0...i_k}(e_{j_0...j_k})
=
\delta^{i_0}{}_{j_0}
...
\delta^{i_k}{}_{j_k}
$$
Any cochain $f\in C^*$ can be written as
$$
f = 
\sum_{k=0}^N 
\sum_{i_0...i_k \in K}
f_{i_0...i_k}
e^{i_0...i_k}
$$
where the $f_{i_0...i_k}$ are scalar coefficients denoting the coordinates of $f$ in the dual basis.
Cochains can be thought as tensor fields defined on top of the simplicial complex $K$.

\subsection{Boundary operator}

The boundary operator $\partial:C\to C$ is a linear operator defined by its action on the basis vectors
$$
\partial
e_{i_0...i_k}
=
\sum_{j=0}^k
(-1)^k
e_{i_0...i_{j-1}i_{j+1}...i_k}
$$
It mimics the behavior of the boundary operator of a ``smooth'' topological space.

\subsection{Coboundary operator}

The coboundary operator $d:C^*\to C^*$ is the linear operator defined by the dual relation
$$
(df)(c)
=
f(\partial c)
$$
for any chain $c\in C$ and any cochain $f\in C^*$.
It mimics the behavior of the differential operator over a ``smooth'' topological space.

The coboundary operator $d$ mimics the role of the gradient operator in ``smooth'' settings.

\subsection{Inner product}

The space of chains can be equipped with an inner product
$$
C^2 \ni c,c' \to \langle c , c'\rangle \in \RR.
$$
For any given and fixed $c\in C$, the mapping
$$
C\ni c' \to \langle c , c'\rangle \in \RR.
$$
defines a covector or, in other words, a cochain denoted by $c^{\flat}$; i.e. an element of $C^*$.
If the inner product is non-singular, the mapping $C\ni c \to c^{\flat} \in C^*$ is invertible.
Its inverse is denoted by
$$
C^* \ni f \to f_{\sharp} \in C.
$$

This invertible mapping together with the inner product on $C$ induces an inner product on $C^*$ given by
$$
C^*\ni f,f' \to \langle f,f' \rangle := \langle f_{\sharp}, f'_{\sharp} \rangle \in \RR
$$
for every $f,f'\in C^*$.

\subsection{Dual boundary operator}

The dual boundary operator $\partial^*:C\to C$ is the inner-product adjoint of the boundary operator $\partial$.
In other words, is defined by
$$
\langle \partial^* c , c' \rangle 
:=
\langle c ,  \partial c' \rangle 
$$
for every $c,c'\in C$.

\subsection{Dual coboundary operator}

The dual coboundary operator $d^*:C^*\to C^*$ is the inner-product adjoint of the coboundary operator $d$.
In other words, is defined by
$$
\langle d^* f , f' \rangle 
:=
\langle f ,  d f' \rangle 
$$
for every $f,f'\in C^*$.

The dual coboundary operator $d^*$ mimics the role of the rotor operator in ``smooth'' settings.

\subsection{Hodge Laplacian}

The Hodge Laplacian is a generalization of the graph Laplacian to simplicial complexes.
It is defined as the linear operator $L:C^*\to C^*$ given by
$$
L = (d+d^*)^2 = dd^* + d^*d
$$
since it can be shown that $d^2 = (d^*)^2 = 0$.

\subsection{Hodge decomposition}


Any cochain $f\in C^*$ admits a so called Hodge decomposition
$$
f
=
g\oplus s\oplus h
$$
where $g,s,h\in C^*$ are unique given $f$.
In particular, the components $g$, $s$ and $h$ are called the gradient, solenoidal and harmonic components of $f$ and are such that:
$$
Lh=0
$$
i.e. $h$ belongs to the kernel of the Hodge Laplacian $L$,
there exists a $w\in C^*$ such that 
$$
g=dw
$$ 
and, there exists a $u\in C^*$ such that 
$$
s=d^*u.
$$

\section{Description of the experiments}

In our numerical experiments we start with an empirically obtained cochain $f$ defined on the links, i.e.
$$
f 
= 
\sum_{ij \in K}
f_{ij} e^{ij}.
$$
Then, we compute its Hodge decomposition to obtain the gradient component
$$
g 
= 
\sum_{ij \in K}
g_{ij} e^{ij},
$$
the solenoidal component
$$
s 
= 
\sum_{ij \in K}
s_{ij} e^{ij}
$$
and the harmonic component
$$
h 
= 
\sum_{ij \in K}
h_{ij} e^{ij}.
$$
Notice, these are also quantities defined over the links (i.e., they look like vector potentials).
The method also provides a ``scalar pontential'' or potential over the nodes $i$:
$$
w
=
\sum_i
w_i e^i
$$
whose ``gradient'' produces the gradient component 
$$
g
=
dw
$$
of $f$.
Crucially, modula an arbitrary constant that is chosen such that $w_i\geq 0$ together with the condition that the smallest rating is $0$, the components $w_i$ provide the ratings of the nodes.
Additionally, the method provides a ``tensor potential'' or potential over the triangles $ijk$:
$$
u
=
\sum_{ijk}
u_{ijk} e^{ijk}
$$
whose ``rotor'' produces the solenoidal component
$$
s
=
d^*
u
$$
of $f$.

\subsection{Discussion}

The ratings $w_i$ in $w$ inferred from $f$ are basically determined by $g$ via the equation $g=dw$ previously mentioned.
Hence, the part $f-g = s\oplus h$ of $f$ that is not explained by $g$ somehow determines how ``inconsistent'' are the inferred ratings.
In particular, $h=0$ for a fully connected network.
As a consequence, we can say that the ratio
$$
\frac{|s|}{|f|} \in [0,1]
$$
somehow quantifies how ``inconsistent'' are the inferred ratings.
If the ratio is zero, the ratings are fully consistent.
If the ratio is one, the ratings make no sense.

%\section{Introducción}

%Complex systems...~\cite{grady2010discrete}.

%\section{Teoría}

%\section{Results}

%See fig.~\ref{fig1}.
%
%\begin{figure*}
%\includegraphics*[scale=.4]{assets/fig1.png}
%\includegraphics*[scale=.4]{assets/fig1.png}
%\put(-493,140){\bf a)}
%\put(-250,140){\bf b)}
%\\
%\includegraphics*[scale=.4]{assets/fig1.png}
%\includegraphics*[scale=.4]{assets/fig1.png}
%\put(-493,140){\bf c)}
%\put(-250,140){\bf d)}
%\vspace{-0.25cm}
%\caption{
%\label{fig1}
%Bla bla.
%{\bf a)} 
%Bla bla.
%{\bf b)} 
%Bla bla.
%{\bf c)} 
%Bla bla.
%{\bf d)} 
%Bla bla.
%}
%\end{figure*}

%\section{Discussion}

%When comparing...

%\section{Conclusions}


% If in two-column mode, this environment will change to single-column
% format so that long equations can be displayed. Use
% sparingly.
%\begin{widetext}
% put long equation here
%\end{widetext}

% figures should be put into the text as floats.
% Use the graphics or graphicx packages (distributed with LaTeX2e)
% and the \includegraphics macro defined in those packages.
% See the LaTeX Graphics Companion by Michel Goosens, Sebastian Rahtz,
% and Frank Mittelbach for instance.
%
% Here is an example of the general form of a figure:
% Fill in the caption in the braces of the \caption{} command. Put the label
% that you will use with \ref{} command in the braces of the \label{} command.
% Use the figure* environment if the figure should span across the
% entire page. There is no need to do explicit centering.

% \begin{figure}
% \includegraphics{}%
% \caption{\label{}}
% \end{figure}

% Surround figure environment with turnpage environment for landscape
% figure
% \begin{turnpage}
% \begin{figure}
% \includegraphics{}%
% \caption{\label{}}
% \end{figure}
% \end{turnpage}

% tables should appear as floats within the text
%
% Here is an example of the general form of a table:
% Fill in the caption in the braces of the \caption{} command. Put the label
% that you will use with \ref{} command in the braces of the \label{} command.
% Insert the column specifiers (l, r, c, d, etc.) in the empty braces of the
% \begin{tabular}{} command.
% The ruledtabular enviroment adds doubled rules to table and sets a
% reasonable default table settings.
% Use the table* environment to get a full-width table in two-column
% Add \usepackage{longtable} and the longtable (or longtable*}
% environment for nicely formatted long tables. Or use the the [H]
% placement option to break a long table (with less control than 
% in longtable).
% \begin{table}%[H] add [H] placement to break table across pages
% \caption{\label{}}
% \begin{ruledtabular}
% \begin{tabular}{}
% Lines of table here ending with \\
% \end{tabular}
% \end{ruledtabular}
% \end{table}

% Surround table environment with turnpage environment for landscape
% table
% \begin{turnpage}
% \begin{table}
% \caption{\label{}}
% \begin{ruledtabular}
% \begin{tabular}{}
% \end{tabular}
% \end{ruledtabular}
% \end{table}
% \end{turnpage}

%\begin{acknowledgments}
%JIP acknowledges financial support from grants CONICET (PIP 112 20150 10028), FonCyT (PICT-2017-0973), SeCyT–UNC (Argentina) and MinCyT C\'ordoba (PID PGC 2018).
%The authors thank CCAD -- Universidad Nacional de C\'ordoba, \url{http://ccad.unc.edu.ar/}, which is part of SNCAD -- MinCyT, Argentina, for the provided computational resources.
%\end{acknowledgments}

% Create the reference section using BibTeX:
\bibliography{ref}

% Specify following sections are appendices. Use \appendix* if there
% only one appendix.

%\onecolumngrid
%\appendix

%\section{Apendice A}


\end{document}
%
% ****** End of file apstemplate.tex ******


